\chapter{Conclusiones y Aportes del Proyecto}
\label{chap:conclusiones}

\par La tesis se propuso transformar un prototipo de reconocimiento de glosa de Lengua de Señas Mexicana en un sistema escalable y sostenible. El logro principal de esta investigación no es solo la precisión del modelo en sí, sino el establecimiento exitoso de un \textbf{ecosistema de MLOps de aprendizaje continuo} que mitiga los desafíos inherentes a la expansión dinámica del léxico.

\par La hipótesis central del proyecto, que postulaba la capacidad de construir un sistema de reconocimiento de LSM robusto y reproducible mediante la implementación de estrategias de aprendizaje continuo y una infraestructura de MLOps, ha sido \textbf{validada empíricamente}.

\section{Aportes Metodológicos y Técnicos Fundamentales}
\label{sec:aportes}

Las contribuciones del proyecto pueden segregarse en dos áreas de impacto directo que dotan de solidez académica al trabajo: la infraestructura de despliegue y la estrategia de Machine Learning.

\subsection{Aporte en MLOps: Reproducibilidad y Despliegue Automatizado}
\par La primera contribución significativa es la transición del desarrollo \textit{offline} (evidenciado por la ejecución inicial de \lstinline|run_training.py|) a un ciclo de \textbf{Integración y Despliegue Continuo para Machine Learning (CD4ML)}. Esto se materializa en:
\begin{enumerate}
    \item \textbf{Orquestación Auditable:} La gestión mediante \lstinline|trainingHandler.js| y la definición de recursos en \lstinline|template.yaml| aseguran que todo el proceso de entrenamiento sea reproducible, versionado y trazable, elementos críticos para la ingeniería de sistemas en producción.
    \item \textbf{Ecosistema Sólido:} El sistema implementado permite que las actualizaciones de datos en S3 se reflejen automáticamente en el entrenamiento incremental (vía \lstinline|model_trainer.py|), separando la lógica de negocio del entorno operativo.
\end{enumerate}

\subsection{Aporte en Aprendizaje Continuo: Mitigación del Olvido Catastrófico}
\par La contribución metodológica más destacada en el ámbito de Machine Learning es la solución eficaz al \textit{olvido catastrófico} (Catastrophic Forgetting) que ocurre al incorporar nuevas señas al léxico. Esta mitigación se logra mediante la \textbf{integración dual} de:
\begin{itemize}
    \item \textit{Weight Anchoring}: Un mecanismo de regularización que prioriza la estabilidad de los pesos neuronales críticos para el conocimiento previamente adquirido.
    \item \textit{Replay Buffer}: La mezcla estratégica de un pequeño conjunto de datos antiguos con los nuevos, garantizando el entrenamiento incremental sin sacrificar la precisión en el léxico base.
\end{itemize}
Este marco establece un precedente para sistemas de reconocimiento de lenguaje de señas que operan en entornos de vocabulario dinámico.

\section{Verificación del Rendimiento Operacional en el Borde}
\par El sistema final valida la eficiencia del modelo \textbf{Bidirectional LSTM} (construido en \lstinline|model_constructor.py|) como una arquitectura apta para la inferencia en dispositivos con recursos limitados (el \textit{Edge}). La aplicación móvil (\lstinline|LiveRecognitionActivity.kt|) funciona como un nodo de inferencia eficiente, utilizando un mecanismo de \textit{voting} que agrega predicciones en el tiempo para incrementar la robustez final en un entorno de streaming de video, demostrando la viabilidad práctica del despliegue en campo.

\section{Limitaciones y Vínculo con el Trabajo Futuro}
\label{sec:limitaciones_futuro}
\par A pesar de los logros mencionados, la madurez académica exige reconocer las limitaciones del proyecto en su iteración actual. Estas deficiencias son precisamente las que definen las líneas de investigación propuestas para el futuro:
\begin{enumerate}
    \item La configuración de la aplicación móvil es \textbf{estática}, lo que impide el ajuste fino remoto de parámetros sensibles de inferencia (\lstinline|REQUIRED_VOTES|) sin una actualización completa de la aplicación.
    \item La \textbf{capacidad de generalización} del modelo se limita a \textit{keypoints} de pose y manos (195 componentes, como se define en \lstinline|constants.py|), omitiendo los cruciales \textit{marcadores no manuales} (MNM) faciales.
    \item La arquitectura BiLSTM, aunque eficiente, es subóptima para capturar dependencias temporales complejas en comparación con modelos de \textit{Self-Attention}.
\end{enumerate}

% =================================================================
% CAPÍTULO DE TRABAJO FUTURO Y LÍNEAS DE INVESTIGACIÓN (Su Contenido Original)
% =================================================================

\section{Trabajo Futuro y Líneas de Investigación}


\par Esta sección resume brevemente las bases del proyecto que justifican las líneas de investigación futuras.

\subsection{Del Prototipo de Entrenamiento al Sistema de MLOps Sólido}

\par La evolución del proyecto ha trascendido significativamente el alcance inicial. Lo que comenzó como un prototipo de entrenamiento \textit{offline} (evidenciado por la primera iteración de la secuencia de \lstinline|run_training.py|) se ha transformado en un sistema de \textbf{MLOps} robusto y reproducible. Este cambio de paradigma se centra en la capacidad de \textbf{entrenamiento incremental} (vía \lstinline|model_trainer.py| y el uso del \textit{Replay Buffer}) y el despliegue orquestado en la nube (gestión mediante \lstinline|trainingHandler.js| y \lstinline|template.yaml|).
\par La principal contribución metodológica de esta evolución radica en la mitigación del \textit{olvido catastrófico} en la expansión del léxico. La integración de técnicas como \textit{Weight Anchoring} y la mezcla de datos antiguos (\textit{Replay}) no solo ha mejorado la precisión, sino que ha establecido un marco sostenible para la incorporación continua de nuevas señas, lo cual es fundamental para cualquier sistema de reconocimiento de lenguaje en constante crecimiento.

\section{Trabajo Futuro I: Hacia la Configuración y Optimización del Borde (\textit{Edge})}
\label{sec:optimizacion_app}

\par Una de las líneas de investigación más prometedoras, y un objetivo central del futuro del proyecto, es trasladar el control de la optimización y la configuración directamente a la capa de la aplicación móvil, buscando la \textbf{optimidad y adaptabilidad} en el punto de inferencia. Esto implica transformar la aplicación de una mera consumidora de un modelo estático a un \textbf{nodo inteligente} dentro del ecosistema de MLOps.

\subsection{Actualizaciones Dinámicas \textit{Over-the-Air} (OTA)}
\par Es imperativo desarrollar un mecanismo para el despliegue de modelos actualizados sin requerir una reinstalación completa de la aplicación. Se propone la implementación de un \textit{endpoint} que permita a la \lstinline|LiveRecognitionActivity.kt| comprobar la existencia de una nueva versión del modelo \lstinline|.tflite| o \lstinline|.json| y descargarla e inicializarla de forma segura en segundo plano. Esto reduciría la latencia de despliegue de nuevos modelos de semanas a minutos.

\subsection{Configuración de Parámetros de Inferencia Centrada en el Usuario}
\par La robustez de la predicción en la aplicación (ver el mecanismo de \textit{voting} en \lstinline|LiveRecognitionActivity.kt|) depende de parámetros sensibles (e.g., el número de \textit{frames} requeridos para el voto, el umbral de confianza). El trabajo futuro debe enfocarse en:
\begin{enumerate}
    \item \textbf{Exposición de Constantes:} Centralizar las constantes críticas de la \textit{app} (ej., \lstinline|REQUIRED_VOTES|, umbrales) para que puedan ser modificadas de forma remota, permitiendo el \textbf{ajuste fino del rendimiento} por entorno o por grupo de usuarios (A/B testing).
    \item \textbf{Optimización del Procesamiento de Keypoints:} La inferencia en el borde exige la máxima eficiencia. Se debe investigar la posibilidad de optimizar el cálculo de la normalización de los \textit{keypoints} (actualmente en \lstinline|data_processor.py|) para ser ejecutado con mínima carga en el hilo de la cámara de Kotlin, explorando bibliotecas de cómputo optimizado para Android.
\end{enumerate}

\section{Trabajo Futuro II: Incremento de la Robustez y Generalización del Modelo}
\label{sec:robustez_modelo}

\par Más allá de la eficiencia operativa, la solidez académica requiere una mejora continua en el poder predictivo y la capacidad de generalización del modelo.

\subsection{Exploración de Arquitecturas Avanzadas (Transformers)}
\par Aunque la arquitectura \textbf{Bidirectional LSTM} ha demostrado ser eficiente (como se observa en \lstinline|model_constructor.py|), el futuro de la clasificación de secuencias dinámicas se inclina hacia los modelos \textbf{Transformer}. Se propone una línea de investigación que reemplace o complemente la capa recurrente por un mecanismo de \textit{Self-Attention}. Esto podría capturar dependencias a largo plazo de manera más efectiva, mejorando la robustez en señas con variaciones temporales significativas.

\subsection{Estrategias Avanzadas de Aprendizaje Incremental}
\par Se propone fortalecer el marco de entrenamiento con técnicas de \textbf{Aprendizaje Continuo} más sofisticadas que el \textit{Replay Buffer} simple, tales como:
\begin{itemize}
    \item \textbf{Fine-Tuning Personalizado (Federated Learning Conceptos):} Investigar métodos de ajuste fino específicos para nuevos usuarios o entornos (ej., diferentes iluminaciones), aunque sea de forma simulada en el \textit{backend}.
    \item \textbf{Regularización del Conocimiento (Knowledge Distillation):} Usar modelos antiguos como "profesores" para guiar el entrenamiento de modelos nuevos, reduciendo el impacto en las clases previamente aprendidas.
\end{itemize}

\subsection{Expansión del \textit{Feature Set}}
\par El vector de \textit{features} actual (195 componentes, véase \lstinline|constants.py|) se basa en \textit{keypoints} de pose y manos. Una mejora crítica sería la incorporación de la \textbf{información facial} (cejas, boca) para capturar los \textit{marcadores no manuales} (MNM). Esta expansión dotaría al modelo de una \textbf{mayor robustez empírica} y una representación lingüística más completa.

\par Estas líneas de trabajo futuras asegurarán que el proyecto mantenga su vanguardia, tanto en la implementación de un sistema de MLOps eficiente como en la mejora de la precisión y generalización del modelo de reconocimiento de LSM.